\documentclass{ForPPGCPR-resumo}
\usepackage{geometry}
\usepackage{enumitem}
\usepackage{lipsum}
\usepackage{expl3}
\usepackage{xcolor}
\usepackage[fontsize=9pt]{fontsize}
\geometry{
 a4paper,
 left=1cm,
 right=1cm,bottom=1cm,top=1.5cm,
 }

\author{
  \fcolorbox{colorForPPGCPR}{colorForPPGCPR}{\color{white}\textbf{P}}rimeiro Autor\\
  \textit{autor1@utfpr.edu.br}
  \and
  \fcolorbox{colorForPPGCPR}{colorForPPGCPR}{\color{white}\textbf{S}}egundo Autor\\
  \textit{autor2@utfpr.edu.br}
  \and
  \fcolorbox{colorForPPGCPR}{colorForPPGCPR}{\color{white}\textbf{T}}erceiro Autor\\
  \textit{autor3@utfpr.edu.br}
}

\title{Título}


\begin{document}
   
\maketitle

\section*{Área de Conhecimento}
\textcolor{colorForPPGCPR}{Indique aqui a área de conhecimento do seu projeto. São exemplos de área de conhecimento: Banco de Dados, Engenharia de Software, Processamento Gráfico, etc. Para mais opções, consultar 
\url{http://lattes.cnpq.br/documents/11871/24930/TabeladeAreasdoConhecimento.pdf/d192ff6b-3e0a-4074-a74d-c280521bd5f7}}

\section*{Objetivo}

        \textcolor{colorForPPGCPR}{Liste aqui o objetivo geral da pesquisa de forma direta e clara.}

\section*{Destaques}
        
\textcolor{colorForPPGCPR}{Nesta parte, adicione de 3 a 5 destaques (\textit{highlights}) da sua pesquisa. Os destaques são uma pequena coleção de pontos principais que capturam os resultados da pesquisa, bem como métodos, técnicas, ferramentas, algoritmos, etc de destaque utilizados.}        
        
\textcolor{colorForPPGCPR}{Veja alguns exemplos em: \url{https://www.elsevier.com/authors/tools-and-resources/highlights}}
        
        \begin{itemize}[leftmargin=*]
	\setlength\itemsep{-0.3em}
	\item \textbf{[1] [}\textcolor{colorForPPGCPR}{Descrição do destaque 1 ...}\textbf{]}
	\item \textbf{[2] [}\textcolor{colorForPPGCPR}{Descrição do destaque 2 ...}\textbf{]}
	\item \textbf{[3] [}\textcolor{colorForPPGCPR}{Descrição do destaque 3 ...}\textbf{]}
	\item \textbf{[4] [}\textcolor{colorForPPGCPR}{Descrição do destaque 4 ...}\textbf{]}
	\item \textbf{[5] [}\textcolor{colorForPPGCPR}{Descrição do destaque 5 ...}\textbf{]}
\end{itemize}

\section*{Resumo}

\textcolor{colorForPPGCPR}{Para elaboração e apresentação do resumo, escolha uma das duas opções a seguir.}

\textcolor{colorForPPGCPR}{
\textbf{[Opção 1: Resumo Gráfico]} Crie uma imagem para ajudar o leitor a obter rapidamente uma visão geral sobre seu trabalho e entender o propósito e os resultados de sua pesquisa. Consultar \url{https://www.elsevier.com/authors/tools-and-resources/visual-abstract} para exemplos e diretrizes para elaboração de bons resumos gráficos.
}

\textcolor{colorForPPGCPR}{\textbf{[Opção 2: Resumo Textual]} Um resumo da pesquisa apresentado de forma textual. }

\section*{Resultados Esperados}
\textcolor{colorForPPGCPR}{Identifique os resultados esperados de sua pesquisa. Por exemplo, no contexto de inovação: são exemplos: produto de software; produto de processo; produto em serviço; etc. Para outras classificações de inovação, consultar \url{https://www.gov.br/capes/pt-br/centrais-de-conteudo/2020-01-03-relatorio-gt-inovacao-e-transferencia-de-conhecimento-pdf}}

\textcolor{colorForPPGCPR}{Outra possibilidade seria na forma de produção técnica. São exemplos: software/aplicativo/programa de computador; tecnologia social; patente; material didático; etc. Para definição de outros tipos de produção técnica, consultar \url{https://www.gov.br/capes/pt-br/centrais-de-conteudo/10062019-producao-tecnica-pdf} – Tabela 2, página 17}



\section*{Trabalhos Relacionados -- Leia também}

\textcolor{colorForPPGCPR}{[Seção opcional] Se aplicável, cite outras referências relacionadas à sua pesquisa, por exemplo, um artigo publicado, uma produção técnica, etc.}

\end{document}